\chapter{System Design (Detailed)}
\label{ch:system}

\section{Overview}

The pipeline maps an OPG to a per-tooth record in six stages. Figure~\ref{fig:flow} draws it in the
style of the Phase 2 target pipeline. Figure~\ref{fig:arch} gives the as-shipped version with each
default checkpoint and its learnable parameter count. Figure~\ref{fig:walkthrough} follows one
held-out radiograph through every stage.

\begin{figure}[!htbp]
\centering
% As-built pipeline in the Phase 2 Figure 3.1 style and palette. Components traced to
% scripts/run_inference.py (docstring stage list, DEF_* checkpoints, CARIES_LESION_CONF 0.25,
% ASSESSABILITY_MIN_FDI_CONFIDENCE 0.30) and src/classifiers.py (crop conventions).
\begin{tikzpicture}[node distance=4.2mm]
\node[io] (in) {OPG X-ray input\\[-1pt]{\normalfont PNG or JPEG, any size}};
\node[stage,below=of in] (s1) {1. Tooth detection\\[-1pt]{\normalfont YOLOv11l, multi-corpus, conf.\ 0.25}};
\node[stage,below=of s1] (s2) {2. FDI numbering\\[-1pt]{\normalfont ResNet18 + box features, Hungarian one-to-one}};
\node[stage,below=of s2] (s3) {3. Per-tooth processing};
\node[cari,below left=7mm and 6mm of s3] (c1) {3a. Lesion detector\\[-1pt]{\normalfont whole image, conf.\ 0.25}};
\node[cari,below=of c1] (c2) {Lesion $\rightarrow$ tooth\\[-1pt]{\normalfont argmax containment}};
\node[cari,below=of c2] (c3) {3b. Severity\\[-1pt]{\normalfont tight crop, ResNet18}};
\node[cari,below=of c3] (c4) {Caries finding\\[-1pt]{\normalfont none / caries / deep}};
\node[bone,below right=7mm and 6mm of s3] (b1) {Expanded crop\\[-1pt]{\normalfont 1.5$\times$ width, 2.0$\times$ height}};
\node[bone,below=of b1] (b2) {4a. Abnormality flag\\[-1pt]{\normalfont ResNet18, radiograph split}};
\node[bone,below=of b2] (b3) {4b. Severity\\[-1pt]{\normalfont withheld by default}};
\node[bone,below=of b3] (b4) {Abnormality finding\\[-1pt]{\normalfont normal / flagged}};
\coordinate (mid) at ($(c4.south)!0.5!(b4.south)+(0,-6mm)$);
\node[stage,below=12mm of mid] (s5) {5. Fusion and assessability\\[-1pt]{\normalfont FDI confidence $\geq$ 0.30, alignment check}};
\node[rept,below=of s5] (s6) {6. Report generation\\[-1pt]{\normalfont JSON \texttt{opg-report/1.6} + overlay}};
\node[io,below=of s6] (out) {Structured per-tooth report\\[-1pt]{\normalfont with limitations and provenance}};
\draw[arr] (in) -- (s1); \draw[arr] (s1) -- (s2); \draw[arr] (s2) -- (s3);
\draw[arr] (s3.south) -- ++(0,-3mm) -| (c1.north);
\draw[arr] (s3.south) -- ++(0,-3mm) -| (b1.north);
\draw[arr] (c1) -- (c2); \draw[arr] (c2) -- (c3); \draw[arr] (c3) -- (c4);
\draw[arr] (b1) -- (b2); \draw[arr] (b2) -- (b3); \draw[arr] (b3) -- (b4);
\draw[arr] (c4.south) |- ($(s5.north)+(0,5mm)$) -- (s5.north);
\draw[arr] (b4.south) |- ($(s5.north)+(0,5mm)$);
\draw[arr] (s5) -- (s6); \draw[arr] (s6) -- (out);
\node[soft=p2teal,text width=21mm,left=3mm of c2,font=\sffamily\tiny] {retained alternative:\\tight-crop classifier\\(\texttt{-{}-caries-method gate})};
\end{tikzpicture}

\caption[As-built pipeline]{The as-built pipeline, drawn in the style and colours of the Phase 2
target architecture. Teal: the caries branch; maroon: the abnormality branch.}
\label{fig:flow}
\end{figure}

\begin{figure}[!htbp]
\centering
\includegraphics[width=\textwidth]{figures/repo/R06_architecture.pdf}
\caption[The pipeline as it ships, with parameter counts]{The pipeline as it ships, with its default
checkpoints and learnable parameter counts read from the saved weights. The default configuration
(Stages 1, 2, default 3a, 3b, 4a, and the withheld 4b with 11.18\,M) totals 95.34\,M. Blue: identity
stages; dark blue: the default caries path; dashed orange: the retained caries variant, the
comparison arm throughout; green: the abnormality flag. Thresholds are fixed in code; none was tuned
on an evaluation cohort.}
\label{fig:arch}
\end{figure}
% {C002} {C003}

\begin{figure}[!htbp]
\centering
\includegraphics[width=0.94\textwidth]{figures/repo/R02_walkthrough.pdf}
\caption[One held-out radiograph through the shipped pipeline]{One held-out DENTEX radiograph
(\texttt{test\_0}, not used to train any model) through the shipped pipeline. (a) Input, 2,829 by
1,316 pixels; dashed: the region shown in (b)--(f). (b) Stage 1: 30 teeth, with detection
confidences. (c) Stage 2: one FDI code per tooth; ``?'' marks its two unsure codes (24, 42).
(d) Stage 3a: one lesion-detector pass over the whole image, four boxes (confidence 0.35 to 0.89).
(e) Each box goes to the tooth holding most of it: 26 and 46 flagged as caries, 36 and 45 as deep
caries; no Stage-4 flag. (f) The written overlay with the expert's three caries annotations (26, 45,
48) in yellow: two found, 48 missed, 36 flagged with no annotated lesion. Radiograph:
DENTEX~\cite{hamamci2023dentex}, CC~BY-NC-SA~4.0.}
\label{fig:walkthrough}
\end{figure}

\section{Stage 1: Tooth Detection}

Stage 1 is a YOLOv11l detector~\cite{jocher2024yolo11} with 25.31\,M learnable parameters. The image is
rotated by 90\textdegree\ counter-clockwise into the orientation the detector was trained on, detection
runs at an input size of 1,024 pixels with confidence threshold 0.25 and non-maximum-suppression IoU
0.50, and the boxes are rotated back into the upright frame. The shipped detector, the default since
11 October 2026, was trained on three corpora (Tufts, DENTEX and the Kaggle dual-labelled set) with
per-corpus balancing. It replaced a Tufts-only YOLOv11l-seg detector, which stays selectable. The
replacement missed one of its pre-registered guards and was adopted by a separate, dated decision for
its numbering and caries behaviour; Section~\ref{sec:res-train} reports both. % {C031}

\section{Stage 2: FDI Numbering}

Stage 2 scores every FDI code for each detected tooth. A crop of the box expanded 1.5 times in each
direction is resized to 224 pixels and passed through a ResNet18~\cite{he2016resnet}, together with four
box-position features (centre $x$, centre $y$, width and height, each divided by the image's full
width or height). The network has two heads, one for the quadrant (four classes) and one for the
position (eight classes). Their log-probabilities are summed into a joint score for each of the 32
codes. The Hungarian algorithm~\cite{kuhn1955hungarian} then chooses the one-to-one matching of teeth to
codes with the highest total joint log-probability, so no code can be emitted twice, which an
independent per-tooth argmax would allow when two adjacent molars both look most like tooth 17. The
reported confidence is the exponential of the chosen code's joint log-probability. The head has
11.19\,M parameters and was trained on DENTEX and Tufts crops.

An optional structural post-pass that enforces the anatomical order of codes along the arch was built
and evaluated under its own pre-registration; it was a null result and ships off by default.

\section{Stage 3: Caries}

\subsection{Stage 3a, Default: Whole-Image Lesion Detector}

The default caries stage is a second YOLOv11l, trained with a single class to draw boxes around
carious lesions anywhere in the radiograph. It needs no tooth identity. Each predicted lesion box is
assigned to the tooth box that contains the largest share of it (intersection area divided by
lesion-box area, ``argmax containment''), and that tooth is flagged as carious. The confidence
threshold, 0.25, was frozen by the pre-registration and never tuned on a cohort. In the report, the
\texttt{gate\_confidence} of a lesion-flagged tooth is the detection score and is marked as such,
because it is not a calibrated probability.

\subsection{Stage 3a, Retained: Per-Crop Classifier}

The retained variant, selected with \texttt{-{}-caries-method gate}, is the Phase 2 design: a ResNet18
judges each tight tooth crop (no expansion, 224 pixels) as carious or not. It was retrained in Phase 3
with a real negative class of healthy teeth. It is kept so that every number measured before the
change of default can be reproduced, and it is the comparison arm of Chapter~\ref{ch:results}.

\subsection{Stage 3b: Severity}

Teeth flagged by either Stage-3a method are graded as caries or deep caries, the two caries labels
DENTEX releases, by the second stage of the Phase 2 cascade: a ResNet18 on the tight crop, trained in
two steps (the final layer first, then the whole network at a low learning rate) with class weights.
It was validated on 133 crops during development and is not evaluated separately in this report.

\section{Stage 4: Abnormality Flag}
\label{sec:sys-stage4}

\subsection{From the Phase 2 Bone-Loss Branch}

Phase 2 built this branch on the Tufts data, and its construction is kept:
\begin{itemize}
\item Annotation parsing joins the expert descriptions, the tooth boxes and the pathology masks on
  the image identifier; box corners are converted from the raw corner format to $[x, y, w, h]$.
\item Each tooth box is expanded asymmetrically, 1.5 times in width and 2.0 times in height, so that
  the crop contains the root and the surrounding bone.
\item A crop is labelled positive when the share of mask pixels inside it exceeds 0.01, an overlap
  ratio rather than an intersection over union, because the masks are sparse relative to the crop.
\end{itemize}
% {C082}

Phase 3 found that the masks behind this label mark expert abnormalities (apical lesions, sclerosis,
cysts), not periodontal bone loss. The field keeps its name, \texttt{bone\_loss}, so that report
consumers do not break, but it is described as an abnormality-proximity flag everywhere, including in
the report's own limitations block.

\subsection{Stage 4a in Phase 3}

The Phase 2 binary model's validation split shared radiographs with its training set
(Section~\ref{sec:res-stage4}). The shipped model was therefore retrained with the split made by
radiograph: a ResNet18 on 224-pixel greyscale crops, focal loss, trained on 678 Tufts radiographs, with
the checkpoint chosen on validation macro-F1 and tested once on 143 held-out radiographs. Three seeds
of the recipe were trained; seed 42 is shipped.

\subsection{Stage 4b: Severity, Withheld}

The Phase 2 severity head grades flagged teeth as mild, moderate or severe. Phase 3 found that its
classes had been derived from crop contrast rather than from the mask overlap they were meant to
encode. Retrained on true overlap bands, it did not beat an image-blind tooth-position prior, so the
grade is withheld from the report unless \texttt{-{}-bone-severity-grade} is given.

\subsection{CEJ-Based Estimation (Phase 2, Exploratory)}

The Phase 2 report also explored a cemento-enamel junction (CEJ) heuristic. Each crop was enhanced
with contrast-limited adaptive histogram equalisation (CLAHE) and a candidate CEJ and bone crest were located, so that bone loss could be expressed as a
CEJ-to-crest distance, the measurement closest to clinical practice. On 296 crops the heuristic was
judged reliable on 62.5\%, and its severity agreed with the overlap labels exactly on 15.1\% and within
one class on 61.6\%. It is kept as a documented negative result and is not part of the pipeline. % {C081}

\section{Stage 5: Fusion and Assessability}

Stage 5 joins the per-stage outputs by box index into one record per tooth. Before joining it checks
that every stage returned one entry per box. A tooth is marked not assessable, with the reason, if it
has no FDI code or its numbering confidence is below 0.30. The records are sorted by FDI code, and the
report summary counts detected and assessable teeth, distinct and duplicate codes, and positive
findings.

\section{Stage 6: Report and Overlay}

Stage 6 writes the JSON report described in Section~\ref{sec:lld-report} and, if asked, an overlay. The
overlay draws each tooth box with its FDI code. Colour encodes the finding: orange for a finding, blue
for none and grey for not assessable. Grey replaced red because red read as the most urgent tooth when
it meant that the pipeline could not say; orange and blue stay distinguishable under red--green colour
blindness. A thin line marks a numbering doubt.

\section{Explainability}

Phase 2 planned Grad-CAM saliency maps for every tooth-level prediction. Phase 3 did not implement
them. The visual evidence the system gives is the overlay: every finding is drawn on the tooth it
refers to, with the lesion boxes that produced it, so a reader can check the claim against the image.
Saliency maps remain future work (Chapter~\ref{ch:conclusion}).
